\documentclass[10pt,a4paper]{article}
\usepackage[margin=2.0cm]{geometry}
\usepackage{booktabs}
\usepackage{multirow}
\usepackage{amsmath}
\usepackage{graphicx}
\usepackage{hyperref}
\usepackage{xcolor}
\usepackage{array}
\hypersetup{colorlinks=true, linkcolor=blue, urlcolor=blue, citecolor=blue}
\setlength{\parskip}{0.28em}
\setlength{\parindent}{0pt}

\title{A Structured Questioner for Collaborative Instance Navigation:\\
Prompt Design, Data Construction, and Model Selection\\
for the CoIN Challenge 2026}
\author{Ning Yang}
\date{August 2026}

\begin{document}
\maketitle

\begin{abstract}
We present a questioner for the CoIN Challenge 2026. The agent is
Qwen3-VL-32B-Instruct, fine-tuned with two-stage QLoRA on (i)~synthetic
instance pairs labeled under a frozen structured prompt and
(ii)~ground-truth-validated pairs from the official training episodes.
A category-only question-deduplication rule is applied at inference.
Development follows a pre-registered ranking of full-success rate (FR),
then success rate (SR), then questions per observation (NQ). On a
leakage-free holdout of 47 official episodes, mixing official data into
Stage~2 (Mix-FT) matches the synthetic-only fine-tune on FR
(0.713 vs.\ 0.709) while cutting NQ from 3.45 to 0.67. The submitted
weights (Full-FT) additionally absorb the holdout episodes with the
same recipe. Ablations cover oracle identity, prompt family, backbone
choice (including a full Qwen3.6-27B replica that does \emph{not}
transfer), ask/conclude ratio, and a rejected chain-rebalancing branch.
\end{abstract}

\section{Task, Metrics, and Protocol}

Given a natural-language description of a target object and a stream of
candidate images, the agent must accept the candidate as the target
instance (score 2), reject it (score 0), or ask the oracle one
clarifying question (score 1)~\cite{zorzi2026coinqa,taioli2025coin}.
Official training data comprise 167 episodes (50 categories, 367
images) and six description types: category, color, context,
color+feature, color+context, and color+context+feature.

\paragraph{Metrics.}
SR is the per-episode fraction of correct decisions, averaged over
episodes: the denominator is the number of successes if the episode is
fully solved, otherwise successes$+1$ (the episode stops at the first
error). FR is the fraction of episodes in which every decision is
correct. NQ/obs is the per-episode ratio of questions to observations,
then averaged. Ask rate is the fraction of episodes with at least one
question. We do \emph{not} use a global $\sum$successes$/\sum$observations
ratio: that statistic weights long episodes and can exceed 1 when
observations are de-duplicated by pixel identity.

\paragraph{Selection rule.}
FR first, SR second, lower NQ as tie-break. All reported numbers use
temperature 0. Unless noted, development uses a frozen local oracle,
Qwen3-VL-32B-Instruct. Model selection never trains on a 47-episode
holdout obtained by splitting official episodes along shared-image
connected components (zero image leakage).

\paragraph{Naming.}
For readability we use paper names rather than checkpoint identifiers.
The frozen inference/data prompt is the \emph{Structured Attribute
Prompt} (SAP). Fine-tuned 32B variants are Syn-FT (synthetic Stage~2
only), Mix-FT (Stage~2 continued on official train-split pairs),
Chain-FT (a rejected rebalancing branch), and Full-FT (submitted;
Mix-FT recipe with holdout pairs included). Release adapters in the
package are labeled \texttt{v3d}/\texttt{v3o}/\texttt{v3p}/\texttt{v3f}
respectively.

\section{Oracle Identity}
\label{sec:oracle}

With the challenge-provided 7B questioner and the paper prompt under
raw decoding, we compared Qwen2.5-VL-32B-Instruct and
Qwen3-VL-32B-Instruct as oracles on all 167 episodes
(Table~\ref{tab:oracle}). Mean $|\Delta\mathrm{SR}|$ is below 0.01 on
every type. Under this low-ask regime the questioner almost never
queries the oracle, so oracle identity cannot move the metric. Qwen3 is
slightly better on color/context and slightly worse on category; we
freeze Qwen3-VL-32B as the development oracle thereafter. The practical
implication is that oracle quality matters mainly in high-ask policies
and in \emph{labeling} (Section~\ref{sec:data}), not in a silent
baseline.

\begin{table}[h]
\centering
\small
\begin{tabular}{lcccc}
\toprule
Type & Qwen2.5-32B SR/FR/NQ & Qwen3-32B SR/FR/NQ & $\Delta$SR \\
\midrule
category & 0.280/0.263/0.009 & 0.274/0.257/0.009 & $-0.006$ \\
color & 0.591/0.485/0.058 & 0.594/0.497/0.068 & $+0.003$ \\
context & 0.567/0.439/0.156 & 0.574/0.445/0.146 & $+0.007$ \\
color+feature & 0.712/0.571/0.138 & 0.707/0.574/0.155 & $-0.005$ \\
color+context & 0.736/0.634/0.201 & 0.741/0.638/0.198 & $+0.004$ \\
color+ctx+feat. & 0.712/0.586/0.267 & 0.706/0.589/0.249 & $-0.005$ \\
\bottomrule
\end{tabular}
\caption{Oracle swap with a low-ask paper prompt and the 7B challenge
questioner, full 167 episodes. Differences are negligible.}
\label{tab:oracle}
\end{table}

\section{Prompt Selection}
\label{sec:prompt}

All prompts share a three-field output schema: a
\texttt{<motivation>} span, a \texttt{<score>} in $\{0,1,2\}$, and a
\texttt{<question>} that is either a single query or \texttt{None}.
We iterate the \emph{instruction} while holding the schema fixed, so
that supervised labels and inference stay aligned.

\subsection{Paper family on the 7B questioner}
Table~\ref{tab:prompt7b} averages six types (full 167 episodes, Qwen3
oracle, raw decoding). \emph{Paper} is the QAsk-Nav wording.
\emph{Paper-YN} adds a yes/no discriminative-question constraint.
\emph{Adaptive} writes an information-dependent ask/conclude rule into
a single prompt (ask when the description is underspecified; conclude
when attributes already suffice). \emph{Uncertainty gating} forces an
ask when the model emits a confident score but hesitant wording.

\begin{table}[h]
\centering
\small
\begin{tabular}{lccc}
\toprule
Prompt / policy & SR & FR & NQ/obs \\
\midrule
Paper + raw decoding & 0.600 & 0.500 & 0.138 \\
Paper-YN + raw decoding & 0.616 & 0.491 & 0.088 \\
Adaptive + raw decoding & 0.615 & 0.491 & 0.089 \\
Paper + uncertainty gating & 0.719 & 0.562 & 0.838 \\
\bottomrule
\end{tabular}
\caption{Prompt and policy means on the 7B questioner (six types,
167 episodes). Uncertainty gating is discarded as a spurious gain.}
\label{tab:prompt7b}
\end{table}

Uncertainty gating looks best on SR/FR but inflates NQ six-fold.
About 70\% of its questions are a canned template (``Besides what is
already stated in \ldots''). The SR lift comes from forcing a second
look at the image, not from informative questions. We discard it.

Paper-YN and Adaptive are essentially tied and both beat Paper.
Adaptive matches Paper-YN without a per-type switch: category/context
SR rise ($+0.029$ / $+0.040$ vs.\ Paper) while NQ on
attribute-rich types falls (e.g.\ color $0.068\rightarrow 0.042$).
Questions under Adaptive on color+feature are genuine yes/no checks,
not templates. We therefore take Adaptive as the best
\emph{single-prompt} member of the paper family.

\subsection{Why a new prompt is needed}
Two failures of Adaptive motivate a rewrite. First, ``the description
is already specific, so do not ask'' treats text informativeness as a
sufficient statistic for the decision. A specific attribute that is
occluded in the current view still requires a question. Second,
discriminative questions anchored on the \emph{candidate} image do not
transfer to the next view; the official protocol instead asks the
oracle about the hidden \emph{target} image in order to build a
reusable instance profile~\cite{zorzi2026coinqa}.

\subsection{Profile prompts, then SAP}
\emph{Profile-V1} forces OBSERVE $\rightarrow$ COMPARE $\rightarrow$
DECIDE in \texttt{<motivation>}, scores each stated attribute as
MATCHED / CONTRADICTED / UNVERIFIABLE, and requires questions to
request a new target-side attribute that is also checkable on
candidates. Putting Profile-V1 on the 7B fine-tune \emph{without}
relabeling fails: the model ignores the three-step format, over-asks
on specific types, and on long contexts degenerates into repeated
tokens. The same prompt on \emph{native} Qwen3-VL-32B succeeds
(Table~\ref{tab:native32}): category SR $0.340\rightarrow 0.436$,
color+feature $0.726\rightarrow 0.753$, with explicit Step~1/2/3
throughout and no loops. The 7B failure was train/test prompt
mismatch, not a bad prompt.

\emph{Profile-V2} compresses reasoning ($<120$ words) and forbids
score~2 unless a distinguishing attribute is present. On native 32B it
is uniformly worse than V1 (mean SR $0.694$ vs.\ $0.709$, FR $0.608$
vs.\ $0.633$) except for a small category SR bump paid for with much
higher NQ. Compactness alone is not an improvement.

An intermediate 32B fine-tune on Profile-V2 labels raised category SR
from native V1's 0.436 to 0.529 but lagged on several specific types.
That split---gains concentrated on type-only descriptions, losses on
richer ones---motivated rebuilding labels under a stricter protocol
rather than shipping the V2 model.

SAP (frozen for all later data and inference) keeps V1's three-step
compare protocol and V2's length cap, and adds stricter decision
rules: score~0 if the described type is absent or any stated attribute
is contradicted; score~1 if the description is type-only or any stated
attribute is unverifiable; score~2 only when every \emph{stated}
attribute matches and none is unverifiable. Questions ask exactly one
new target attribute. On native 32B, SAP raises mean SR/FR to
$0.739/0.653$ with NQ $0.402$ (category NQ $1.11$, ask rate $99\%$).
It is not uniformly best versus V1 on every attribute-rich type (V1
files there are slightly incomplete, $n=127$--$164$), but it is the
protocol we need in \emph{labels}: every stated attribute is audited,
type-only descriptions must ask, and the same text is used at train
and test. We freeze SAP from data generation onward.

\begin{table}[h]
\centering
\small
\begin{tabular}{lcccc}
\toprule
Setting (native 32B, raw decoding) & SR & FR & NQ/obs & Ask rate \\
\midrule
Adaptive, category & 0.340 & 0.311 & 0.162 & 22\% \\
Profile-V1, category & \textbf{0.436} & \textbf{0.395} & 0.436 & 49\% \\
Adaptive, color+feature & 0.726 & 0.640 & 0.030 & 7\% \\
Profile-V1, color+feature & \textbf{0.753} & \textbf{0.659} & 0.077 & 9\% \\
Profile-V1 (six-type mean$^\dagger$) & \textbf{0.709} & \textbf{0.633} & 0.128 & 16\% \\
Profile-V2 (six types, $n{=}167$) & 0.694 & 0.608 & 0.164 & 20\% \\
SAP (six types, $n{=}167$) & 0.739 & 0.653 & 0.402 & 40\% \\
\bottomrule
\end{tabular}
\caption{Native Qwen3-VL-32B prompt comparison. $^\dagger$Some V1
attribute-type files are slightly incomplete ($n=127$--$164$); the mean
is the best available V1 estimate.}
\label{tab:native32}
\end{table}

\section{Backbone Selection}
\label{sec:backbone}

Replacing the 7B challenge questioner with native Qwen3-VL-32B under
SAP is the first large jump (Section~\ref{sec:prompt}). We then swap
only the questioner to native Qwen3.6-27B, keeping SAP, raw decoding,
temperature 0, and the same 32B oracle (Table~\ref{tab:q36zs}). Mean
SR/FR rise ($0.739\rightarrow 0.748$, $0.653\rightarrow 0.665$) and
NQ drops 55\% ($0.402\rightarrow 0.182$). Five of six types improve;
only context regresses. Category NQ falls from $1.11$ to $0.70$.
(Category $n{=}166$ for Qwen3.6, one missing stove episode; other
types $n{=}167$.) Native Qwen3.6 is the stronger \emph{zero-shot}
questioner.

We nevertheless fine-tune Qwen3-VL-32B: the QLoRA, image-size, and
serving stack were validated on that architecture, and a stronger
zero-shot prior does not imply that an ask-heavy supervised recipe
will transfer. Section~\ref{sec:q36ft} reports a full replica that
loses on holdout FR; the submitted backbone remains 32B.

\begin{table}[h]
\centering
\small
\begin{tabular}{lcc}
\toprule
Type & Qwen3.6-27B SR/FR/NQ & Qwen3-VL-32B SR/FR/NQ \\
\midrule
category & 0.679/0.590/0.699 & 0.659/0.575/1.114 \\
color & 0.682/0.599/0.066 & 0.672/0.557/0.295 \\
context & 0.648/0.557/0.146 & 0.685/0.605/0.351 \\
color+feature & 0.768/0.683/0.046 & 0.742/0.653/0.351 \\
color+context & 0.845/0.767/0.068 & 0.836/0.749/0.172 \\
color+ctx+feat. & 0.868/0.796/0.071 & 0.839/0.778/0.127 \\
\midrule
Mean & \textbf{0.748}/\textbf{0.665}/\textbf{0.182} & 0.739/0.653/0.402 \\
\bottomrule
\end{tabular}
\caption{Zero-shot SAP, same oracle and policy, full official set.
Qwen3.6 asks less and scores slightly higher.}
\label{tab:q36zs}
\end{table}

\section{Data Construction}
\label{sec:data}

Labels are generated with SAP frozen, so the supervised distribution
matches inference. We never train on raw code-variable dumps; each
sample is a (description, candidate image, optional dialogue,
motivation, score, question) tuple.

\subsection{Synthetic instances}
Candidate images come primarily from Habitat renders of CoIN-Bench /
HM3D episodes (target views and nearby viewpoints), plus
same-category cross-instance pairing as hard negatives. InstructPix2Pix
appearance edits were prototyped and then demoted: geometric renders
are the main source. Each pair is labeled by Gemini Flash under SAP.
Strata with low Flash/Pro agreement are escalated to Gemini Pro (Pro
is the final arbiter); invalid or persistently conflicting items are
quarantined rather than force-labeled.

Multi-turn chains (0--6 rounds) are collected with Flash as questioner
and Pro as answerer against the target reference. We then delete
degenerate behavior: near-duplicate questions, answers that contradict
already-stated history, and reasoning that ignores available context
(the dominant failure mode of an early fine-tune). The
train/validation cut is at instance level, not image crop level.

\subsection{Two-stage supervision}
Following QAsk-Nav~\cite{zorzi2026coinqa}: Stage~1 trains motivation
plus score on all labeled pairs (20{,}653 examples) with empty
question slots, to stabilize the schema. Stage~2 continues on the
full schema with an ask-weighted mix. A 10:1 ask:conclude branch is
dominated by 5:1 on a 60-episode slice (SR 0.773 / FR 0.669 / NQ 7.12
vs.\ 0.806 / 0.700 / 4.63); extra ask mass amplifies repeat questions.
We keep 5:1 (8{,}259 Stage~2 examples). Training is QLoRA, 4-bit NF4,
LoRA rank 16 ($\approx 40$M trainable parameters on 32B), images
resized to $224^2$, 7-GPU DDP, learning rate $1\times 10^{-4}$ for
Stages~1--2 and $5\times 10^{-5}$ for official continuations, one epoch
per stage.

\subsection{Official episodes, without leakage}
Synthetic images do not overlap the official 367 files. To use official
episodes without contaminating the test of generalization, we split the
167 episodes by shared-image connected components into 120 train / 47
holdout episodes. Holdout images never appear in Mix-FT training.

On the 120-episode split we label all description$\times$candidate
pairs (2{,}298) with Flash under SAP and \emph{validate against episode
ground truth}:
\begin{itemize}
\item agreeing conclusions are kept (1{,}484);
\item Flash-match vs.\ GT-reject conflicts (296 of 302 disagreements)
are dropped from conclusion training---these are exactly the
underspecified cases the benchmark wants the agent to \emph{ask}
about, so teaching a hard accept would be label noise;
\item depth-0 ask outputs are kept (496).
\end{itemize}
This yields 1{,}980 in-domain examples mixed into Stage~2 (10{,}239
total) for a one-epoch continuation from Syn-FT (Mix-FT).

After selecting Mix-FT we apply the same pipeline to the 47 holdout
episodes (870 pairs $\rightarrow$ 716 kept: 190 ask / 229 match / 297
reject) and continue from the same Syn-FT initialization (Full-FT;
10{,}955 Stage~2 examples). Selection numbers always cite Mix-FT on
holdout; Full-FT is the submitted weight.

\subsection{A rejected rebalancing branch}
Mix-FT regresses on context and color+feature, consistent with
official conclusions suppressing questions. We built 636
GT-checked one-round chains (138 forced asks, 355 contextual rejects,
143 matches) and down-sampled official rejects to $1.2\times$ matches
(Chain-FT). Holdout FR fell to 0.677 (Table~\ref{tab:holdout}). At this
scale the natural official label mix is closer to optimal than
hand-rebalanced chains. Chain-FT is not used.

\section{Inference Policy}
\label{sec:policy}

Syn-FT on category-only descriptions enters repeat-ask loops. A
lightweight rule detects near-duplicate questions (Jaccard $\geq 0.6$
over content tokens of length $\geq 2$) or budget exhaustion, then
re-prompts once to decide from answers already collected (fallback:
reject). Applying this rule to \emph{all} types cuts NQ to 0.74 but
costs up to 13pp FR on color. Restricting it to category is strictly
helpful on the full 167 set (category SR $0.663\rightarrow 0.725$,
FR $0.569\rightarrow 0.641$, NQ $10.2\rightarrow 1.05$). All reported
fine-tune numbers use this hybrid policy: forced decision on category,
raw decoding otherwise.

\section{Main Results and Selection}
\label{sec:main}

\subsection{Holdout-47 (clean benchmark)}
Table~\ref{tab:holdout} is the selection table. Syn-FT raises FR over
zero-shot SAP but explodes NQ. Mix-FT keeps FR (0.713 vs.\ 0.709) and
reduces NQ five-fold. Per type, official data helps category
(FR $0.596\rightarrow 0.745$, SR $0.713\rightarrow 0.840$) and
color+context (FR $0.766\rightarrow 0.830$), with some loss on context
and color+feature. Under the pre-registered rule Mix-FT wins. Chain-FT
is dominated on FR.

\begin{table}[h]
\centering
\small
\begin{tabular}{lccc}
\toprule
Qwen3-VL-32B, hybrid policy, holdout-47 & SR & FR & NQ/obs \\
\midrule
Zero-shot + SAP & 0.766 & 0.688 & 0.41 \\
Syn-FT & \textbf{0.823} & 0.709 & 3.45 \\
Mix-FT (selected) & 0.801 & \textbf{0.713} & 0.67 \\
Chain-FT (not used) & 0.772 & 0.677 & 0.54 \\
\bottomrule
\end{tabular}
\caption{Leakage-free holdout, six types averaged. Mix-FT is selected
by FR, then SR, then NQ.}
\label{tab:holdout}
\end{table}

\begin{table}[h]
\centering
\small
\begin{tabular}{lcc}
\toprule
Type & Syn-FT SR/FR/NQ & Mix-FT SR/FR/NQ \\
\midrule
category & 0.713/0.596/0.98 & \textbf{0.840}/\textbf{0.745}/0.83 \\
color & 0.836/0.681/5.58 & 0.779/0.681/\textbf{1.20} \\
context & \textbf{0.881}/\textbf{0.787}/3.21 & 0.720/0.660/0.56 \\
color+feature & 0.817/0.660/3.11 & 0.755/0.596/0.53 \\
color+context & 0.857/0.766/3.72 & 0.872/\textbf{0.830}/0.70 \\
color+ctx+feat. & 0.833/0.766/4.09 & 0.840/0.766/\textbf{0.18} \\
\midrule
Mean & 0.823/0.709/3.45 & 0.801/\textbf{0.713}/\textbf{0.67} \\
\bottomrule
\end{tabular}
\caption{Holdout-47 by description type. Official pairs mainly repair
category behavior and suppress redundant questions.}
\label{tab:holdout-type}
\end{table}

\subsection{Full-FT sanity}
Continuing from Syn-FT with all 167 official episodes included, a
same-slice sanity on the first 60 episodes per type shows no
regression vs.\ Mix-FT: Mix-FT SR/FR/NQ $0.768/0.667/0.68$ vs.\
Full-FT $0.768/0.678/0.61$, with category FR $0.517\rightarrow 0.633$.
Full-FT is the submitted weight. Clean-benchmark numbers remain Mix-FT
on holdout-47.

On the \emph{full} 167 episodes, Syn-FT plus the hybrid policy already
beats native SAP (mean SR $0.814$ / FR $0.714$ / NQ $3.88$ vs.\
$0.739/0.653/0.40$). Mix-FT/Full-FT exist to keep that accuracy while
making NQ competition-plausible.

\subsection{Replica on Qwen3.6-27B}
\label{sec:q36ft}
We repeat Stage~1, Stage~2 (5:1), and Mix-FT on Qwen3.6-27B with the
same files, hyperparameters, holdout, SAP, and hybrid policy
(Table~\ref{tab:q36ft}). Native Qwen3.6 was stronger; after the same
SFT recipe both Syn-FT and Mix-FT trail the 32B counterparts on FR
and SR. Mix-FT NQ is competitive (0.54 vs.\ 0.67) but cannot overturn
the FR gap ($-6.4$pp). We therefore do not switch the submission.
The likely cause is a mismatch between an already-conservative
zero-shot prior and ask-heavy Stage~2 data, not a failed load: a
smoke test of the 3.6 adapter produces well-formed SAP outputs.

\begin{table}[h]
\centering
\small
\begin{tabular}{lcc}
\toprule
Holdout-47, hybrid policy & Qwen3.6-27B & Qwen3-VL-32B \\
\midrule
Syn-FT mean SR/FR/NQ & 0.779/0.642/3.15 & \textbf{0.823}/\textbf{0.709}/3.45 \\
Mix-FT mean SR/FR/NQ & 0.768/0.649/0.54 & \textbf{0.801}/\textbf{0.713}/0.67 \\
\bottomrule
\end{tabular}
\caption{Same data and recipe on two backbones. Fine-tuning does not
preserve Qwen3.6's zero-shot advantage; 32B Mix-FT remains selected.}
\label{tab:q36ft}
\end{table}

\section{Submitted System}

Questioner: Qwen3-VL-32B-Instruct with the Full-FT LoRA adapter
(rank 16), SAP, temperature 0, hybrid policy (forced decision only on
category-only descriptions). Development oracle: Qwen3-VL-32B-Instruct.
The package includes the adapter, optional merged bf16 weights, Mix-FT
as a clean-benchmark backup, data and training scripts, and the frozen
prompt. Serve with vLLM \texttt{--enable-lora} on the bf16 base, or
serve the merged checkpoint directly.

\section{Reproducibility Notes}
Connected-component episode IDs, Stage~2 manifests, and the exact SAP
string live in the code directory. Metrics are recomputed from dumped
JSON with a single script that implements Section~1. Continuation runs
always start from Syn-FT, not from Mix-FT, so Mix-FT and Full-FT are
parallel forks rather than a chain.

\begin{thebibliography}{9}
\bibitem{taioli2025coin}
F.~Taioli, E.~Zorzi, G.~Franchi, A.~Castellini, A.~Farinelli,
M.~Cristani, Y.~Wang.
Collaborative Instance Object Navigation: Leveraging Uncertainty-Awareness
to Minimize Human-Agent Dialogues. In \emph{ICCV}, 2025.

\bibitem{zorzi2026coinqa}
E.~Zorzi, F.~Taioli, Y.~Wang, M.~Cristani, A.~Farinelli, A.~Castellini,
L.~Bazzani.
Benchmarking Interaction, Beyond Policy: a Reproducible Benchmark for
Collaborative Instance Object Navigation. arXiv:2604.00265, 2026.
\end{thebibliography}

\end{document}
